Ce projet de fin d'année a permis de concevoir et de réaliser \textbf{QuickTriage}, un système
intelligent d'aide au triage des patients aux urgences, couvrant l'intégralité de la chaîne, depuis
la donnée brute jusqu'à l'application mobile destinée au personnel soignant.

\medskip
\noindent\textbf{Bilan du travail réalisé.} Le projet a d'abord posé le cadre clinique et la
problématique du triage, puis dressé un état de l'art des échelles normalisées (CTAS, ESI, MTS,
KTAS) et des approches d'apprentissage automatique appliquées au triage. Sur cette base, les besoins
fonctionnels et non-fonctionnels ont été formalisés, puis traduits en une architecture en trois
couches faiblement couplées : un pipeline d'apprentissage en Python, un backend Jakarta EE exposant
le modèle au format ONNX via une API REST JAX-RS, et un frontend mobile Flutter. La modélisation UML
(cas d'utilisation, classes, séquences) a précisé la conception, avant que l'implémentation ne la
concrétise.

\medskip
\noindent\textbf{Principaux résultats.} Le modèle XGBoost final, optimisé par Optuna et entraîné sur
un jeu de données rééquilibré par augmentation hybride, atteint une exactitude de 73,8~\%, un
F1-score pondéré de 0,729 et, surtout, un rappel de 83,9~\% sur la classe \textsc{Critique}. Ce
dernier résultat, obtenu grâce à une calibration du seuil de décision à 0,30, traduit le parti pris
fondamental du projet : \textbf{minimiser le sous-triage} pour préserver la sécurité des patients en
danger vital. La détection et la correction d'une fuite de données ont par ailleurs garanti
l'honnêteté et la généralisabilité des performances annoncées.

\medskip
\noindent\textbf{Limites.} Plusieurs limites doivent être soulignées. Les performances restent
contraintes par la taille et la nature du jeu de données initial (1~267 patients réels) ; le recours
à l'augmentation synthétique, bien que nécessaire, introduit une part d'artificialité dans les
données d'entraînement. De plus, le système n'a pas été validé en conditions cliniques réelles, et
l'historisation comme la sécurisation complète (authentification, HTTPS de bout en bout) relèvent
encore d'une mise en production à venir.

\medskip
\noindent\textbf{Perspectives.} Plusieurs axes d'amélioration se dégagent :
\begin{itemize}
  \item \textbf{Enrichissement des données} : intégrer des jeux de données réels plus volumineux et
  multicentriques, afin de réduire la dépendance à l'augmentation synthétique et d'améliorer la
  généralisation.
  \item \textbf{Validation clinique} : conduire une évaluation prospective auprès de soignants, pour
  mesurer l'apport réel de l'outil et son acceptabilité sur le terrain.
  \item \textbf{Explicabilité} : intégrer des méthodes d'explication (par exemple SHAP) afin de
  rendre transparentes les raisons d'une prédiction et de renforcer la confiance des soignants.
  \item \textbf{Sécurité et conformité} : finaliser l'authentification, le chiffrement de bout en
  bout et la conformité réglementaire (RGPD) pour un déploiement en production.
  \item \textbf{Apprentissage continu} : exploiter les décisions validées et corrigées par les
  médecins (historisées) pour réentraîner périodiquement le modèle et l'améliorer dans le temps.
\end{itemize}

\medskip
\noindent En définitive, QuickTriage démontre la faisabilité d'une chaîne complète, interopérable et
orientée sécurité patient pour l'aide au triage. Au-delà de ses résultats techniques, le projet
illustre une démarche d'ingénierie rigoureuse, attentive aux pièges méthodologiques de la science des
données et à l'impératif éthique de placer le jugement clinique du soignant au cœur de la décision.
